\documentclass[11pt]{article}
\usepackage[utf8]{inputenc}
\usepackage[margin=1in]{geometry}
\usepackage{amsmath}
\usepackage{graphicx}
\usepackage{booktabs}
\usepackage{hyperref}
\usepackage[round]{natbib}
\usepackage{float}

\title{Evaluating the Biome Smartphone App for Biodiversity Monitoring: Data Quality and Impact on Species Distribution Modeling in Japan}
\author{Anonymous Authors}
\date{}

\begin{document}
\maketitle

\begin{abstract}
Comprehensive biodiversity data is essential for conservation planning, yet traditional expert surveys cannot achieve sufficient spatiotemporal resolution alone. Community science platforms accumulate data rapidly but may harbor spatial and taxonomic biases. Here we systematically evaluate the Biome smartphone application, which has accumulated over 5.27 million observations of more than 40,957 species across Japan since 2019. Species identification accuracy exceeds 95\% for birds, reptiles, mammals, and amphibians, with seed plants at approximately 90\% and insects, mollusks, and fishes below 90\%. We built species distribution models (SDMs) for 132 terrestrial plants and animals using MaxEnt, comparing three data configurations: traditional survey data only, Biome data only, and blended. Models incorporating Biome data consistently outperformed traditional-only models, as measured by the Boyce index with spatial block cross-validation. For endangered species, traditional data alone required over 2,000 records for accurate models (Boyce index $\geq 0.9$), while blending reduced this threshold to approximately 300 records. Principal component analysis revealed that Biome data provides uniform coverage across the urban--natural environmental gradient, whereas traditional data is biased toward natural areas ($p < 0.001$). Optimal model accuracy was achieved when training data consisted of 50--70\% Biome records. These findings demonstrate that community-sourced smartphone data can substantially enhance biodiversity monitoring and conservation planning when appropriately blended with traditional surveys.
\end{abstract}

\section{Introduction}

The ongoing biodiversity crisis demands monitoring systems with unprecedented spatiotemporal resolution \citep{hochkirch2021framework}. International frameworks including the Kunming-Montreal Global Biodiversity Framework (30-by-30 targets) and the Taskforce on Nature-related Financial Disclosures (TNFD) require comprehensive species occurrence data for effective conservation planning and corporate accountability \citep{cbd2022framework, tnfd2023recommendations}. Traditional expert surveys, while providing high-quality data, cannot achieve the geographic and temporal coverage needed to meet these demands \citep{chandler2017contribution}.

Community science (citizen science) platforms have emerged as a powerful complement to traditional surveys, leveraging public participation to accumulate vast quantities of species occurrence records \citep{dickinson2010citizen, bonney2009citizen}. Platforms such as iNaturalist and eBird have demonstrated the potential of community-collected data for ecological research \citep{sullivan2009ebird, loarie2016inaturalist}. However, community science data may suffer from spatial biases (concentration in accessible areas near roads and population centers) and taxonomic biases (misidentification of difficult species groups) \citep{beck2014spatial, geldmann2016determines, kosmala2016assessing}.

The Biome smartphone application, launched in April 2019 in Japan, represents a new generation of community science tools that combine CNN-based image recognition with geospatial data and gamification features to incentivize participation \citep{lecun2015deep}. Despite rapid data accumulation, no systematic evaluation of Biome's data quality across taxa and rarity levels had been conducted, and the potential for smartphone-driven community data to complement traditional surveys in species distribution modeling (SDM) had not been quantified.

Here we address three questions: (1) What is the identification accuracy of Biome data across taxonomic groups and species rarity levels? (2) Can blending Biome and traditional data improve SDM accuracy? (3) What is the optimal blending ratio, and how does it affect the minimum sample size needed for reliable models of endangered species?

\section{Results}

\subsection{Data Accumulation and Filtering}

By July 2023, the Biome app had accumulated over 5.27 million observations (Table~\ref{tab:overview}).

\begin{table}[H]
\centering
\caption{Biome application data overview as of July 7, 2023.}
\label{tab:overview}
\begin{tabular}{lr}
\toprule
\textbf{Metric} & \textbf{Value} \\
\midrule
Total occurrence records & 5,275,457 \\
Species documented & 40,957 \\
App downloads (by Sept 2023) & 839,844 \\
Average daily submissions (2022) & 5,407 \\
Records passing automatic filtering & 2,373,303 \\
Filtering pass rate & 45.0\% \\
Launch date & April 2019 \\
\bottomrule
\end{tabular}
\end{table}

The automatic filtering pipeline excluded unclear images, non-wild individuals, and duplicate observations, retaining 45.0\% of total submissions for scientific analysis.

\subsection{Taxonomic Composition}

Seed plants and insects dominated the dataset (Table~\ref{tab:composition}).

\begin{table}[H]
\centering
\caption{Taxonomic composition of filtered Biome records.}
\label{tab:composition}
\begin{tabular}{lrr}
\toprule
\textbf{Taxon} & \textbf{Fraction (\%)} & \textbf{Fraction Wild (\%)} \\
\midrule
Seed plants & 41.8 & $<90$ \\
Insects (incl.\ Arachnids) & 31.2 & $>97$ \\
Birds & 8.4 & $>97$ \\
Mammals & 3.1 & $<90$ \\
Reptiles & 2.7 & $>90$ \\
Amphibians & 2.3 & $>90$ \\
Fishes & 4.2 & $<90$ \\
Mollusks & 3.8 & $<90$ \\
Other & 2.5 & -- \\
\bottomrule
\end{tabular}
\end{table}

\subsection{Species Identification Accuracy}

Expert verification revealed taxon-dependent identification accuracy (Table~\ref{tab:accuracy}).

\begin{table}[H]
\centering
\caption{Species identification accuracy assessed by expert verification of Biome records, stratified by taxonomic group. Accuracy at species, genus, and family levels among wild-individual records.}
\label{tab:accuracy}
\begin{tabular}{lrrr}
\toprule
\textbf{Taxon} & \textbf{Species (\%)} & \textbf{Genus (\%)} & \textbf{Family (\%)} \\
\midrule
Birds & 96.3 & 98.1 & 99.2 \\
Reptiles & 95.8 & 97.4 & 98.7 \\
Mammals & 95.1 & 97.8 & 99.0 \\
Amphibians & 95.4 & 97.2 & 98.9 \\
Seed plants & 89.7 & 93.8 & 96.4 \\
Insects & 82.4 & 87.6 & 93.1 \\
Fishes & 84.2 & 91.3 & 95.8 \\
Mollusks & 81.7 & 90.4 & 94.2 \\
\bottomrule
\end{tabular}
\end{table}

Birds, reptiles, mammals, and amphibians all exceeded 95\% species-level accuracy. Seed plants achieved approximately 90\%, while insects, fishes, and mollusks fell below 90\%. These accuracy levels substantially exceed those reported for other community platforms: a comparative benchmark of PlantNet, PlantSnap, LeafSnap, iNaturalist, and Google Lens achieved only 69\% identification accuracy for plants \citep{joly2016lifeclef}.

\subsection{Environmental Gradient Coverage}

Principal component analysis of environmental variables revealed fundamental differences in spatial coverage between data sources (Table~\ref{tab:pca}).

\begin{table}[H]
\centering
\caption{Environmental coverage analysis. PC1 captures the urban--natural gradient driven by land use, topography, and climate variables.}
\label{tab:pca}
\begin{tabular}{lrr}
\toprule
\textbf{Metric} & \textbf{Biome Data} & \textbf{Traditional Data} \\
\midrule
PC1 variance explained (\%) & \multicolumn{2}{c}{6.1} \\
Coverage of urban gradient & Uniform & Sparse \\
Coverage of natural gradient & Uniform & Dense \\
Bias direction & Minimal & Toward natural areas \\
KS test (Biome vs.\ uniform) & $p = 0.312$ & -- \\
KS test (Traditional vs.\ uniform) & -- & $p < 0.001$ \\
\bottomrule
\end{tabular}
\end{table}

Biome data was indistinguishable from uniform coverage across the urban--natural gradient (Kolmogorov-Smirnov test, $p = 0.312$), while traditional data was significantly biased toward natural areas ($p < 0.001$). This complementary coverage explains why blending the data sources improves SDM accuracy.

\subsection{Species Distribution Model Performance}

SDMs were built for 132 terrestrial species using MaxEnt \citep{phillips2006maximum, elith2011statistical} with spatial block cross-validation \citep{roberts2017cross} and evaluated using the Boyce index \citep{boyce2002evaluating}.

\begin{table}[H]
\centering
\caption{SDM accuracy (Boyce index) by data source and record count threshold. Values are median across 132 species with 3 replicates each.}
\label{tab:sdm}
\begin{tabular}{lrrr}
\toprule
\textbf{Record Count} & \textbf{Traditional Only} & \textbf{Biome + Traditional} & \textbf{Improvement} \\
\midrule
100 & 0.62 & 0.74 & $+0.12$ \\
300 & 0.71 & 0.87 & $+0.16$ \\
500 & 0.78 & 0.91 & $+0.13$ \\
1,000 & 0.84 & 0.93 & $+0.09$ \\
2,000 & 0.90 & 0.95 & $+0.05$ \\
5,000 & 0.93 & 0.96 & $+0.03$ \\
\bottomrule
\end{tabular}
\end{table}

Blended models consistently outperformed traditional-only models across all record count thresholds, with the largest improvements at lower sample sizes. The improvement was most dramatic for endangered species, where the threshold for achieving Boyce index $\geq 0.9$ dropped from over 2,000 to approximately 300 records.

\subsection{Optimal Blending Ratio}

Systematic variation of the Biome data fraction in blended models revealed an optimum at 50--70\% (Table~\ref{tab:blending}).

\begin{table}[H]
\centering
\caption{Median Boyce index across 132 species as a function of Biome data fraction in the blended training dataset.}
\label{tab:blending}
\begin{tabular}{lr}
\toprule
\textbf{Biome Fraction (\%)} & \textbf{Median Boyce Index} \\
\midrule
0 (Traditional only) & 0.84 \\
10 & 0.86 \\
30 & 0.89 \\
50 & \textbf{0.93} \\
70 & \textbf{0.93} \\
90 & 0.88 \\
100 (Biome only) & 0.82 \\
\bottomrule
\end{tabular}
\end{table}

Model accuracy peaked at 50--70\% Biome data, declining at both extremes. Using only traditional data underperforms due to environmental coverage bias, while using only Biome data underperforms due to lower identification accuracy for some taxa.

\section{Discussion}

Our evaluation demonstrates that the Biome smartphone application provides a substantial and largely untapped resource for biodiversity monitoring in Japan. With over 5.27 million observations and species-level identification accuracy exceeding 95\% for vertebrates, Biome data meets quality standards suitable for ecological research \citep{kosmala2016assessing}. The key finding that blended data reduces the sample size threshold for accurate endangered species SDMs from over 2,000 to approximately 300 records has direct conservation implications, as many endangered species have limited occurrence records in traditional databases \citep{callaghan2019improving}.

The complementary environmental coverage of Biome and traditional data sources drives the improvement in SDM accuracy. Traditional surveys, conducted primarily in natural reserves and protected areas, systematically underrepresent urban and peri-urban environments where many species also occur \citep{beck2014spatial}. Biome's gamification features incentivize users to photograph organisms in their daily environments, naturally generating records across the full urban--natural gradient. This addresses a fundamental limitation of traditional survey design for SDM construction \citep{geldmann2016determines}.

The optimal blending ratio of 50--70\% Biome data indicates that neither data source alone is sufficient for maximally accurate models. Traditional data contributes higher identification accuracy (particularly for difficult taxa) and more systematic sampling in remote natural areas, while Biome data provides uniform environmental coverage and rapid data accumulation \citep{pocock2018vision}.

Limitations include the lower identification accuracy for insects, fishes, and mollusks, which require continued improvement in AI identification algorithms and expanded expert review capacity. Additionally, our analysis focused on Japanese species, and generalizability to other regions with different species assemblages and user demographics should be tested. The 45\% filtering pass rate, while ensuring data quality, means that the majority of submitted records are excluded, and improved user training could increase this rate.

\section{Methods}

\subsection{Biome Application}

The Biome app uses CNN-based image recognition combined with geospatial occurrence data to suggest candidate species identifications. Users photograph organisms; the app records GPS coordinates and timestamps from EXIF data. Gamification features including points for rare and endangered species, quests, and a leveling system incentivize participation. An automatic filtering pipeline excludes unclear images, non-wild individuals, and duplicate observations.

\subsection{Data Quality Assessment}

Expert reviewers verified identification accuracy at species, genus, and family levels across taxonomic groups, stratified by species rarity. A random sample of records was drawn from each taxon, and experts with taxonomic specialization assessed correctness.

\subsection{Species Distribution Modeling}

SDMs were built for 132 terrestrial species using MaxEnt \citep{phillips2006maximum} with environmental predictors including climate (temperature, precipitation), land use (urban, agricultural, forest cover), and topography (elevation, slope). Three data configurations were compared: traditional only, Biome only, and blended. Model accuracy was evaluated using the Boyce index \citep{boyce2002evaluating} with spatial block cross-validation \citep{roberts2017cross} and 3 replicates per species per record count.

\subsection{Environmental Gradient Analysis}

PCA was performed on environmental variables to characterize the urban--natural gradient. The distribution of Biome and traditional data points along PC1 was compared using Kolmogorov-Smirnov tests to assess uniformity of environmental coverage.

\section{Conclusion}

The Biome smartphone application represents a valuable complement to traditional biodiversity surveys, providing high-accuracy species occurrence data with uniform environmental coverage. Blending Biome and traditional data substantially improves species distribution model accuracy and reduces sample size requirements for endangered species, with optimal performance at 50--70\% Biome data. These findings support the integration of community science smartphone platforms into national and international biodiversity monitoring frameworks.

\bibliographystyle{plainnat}
\bibliography{references}

\end{document}
